\documentclass[aspectratio=169,12pt]{beamer}
\usepackage{ugd-official-beamer}

\begin{document}
	\title{Наслов на презентацијата}
	\subtitle{Кратка порака}
	\date[1.1.2027]{1 јануари 2027, 11:00 часот}
	\institute{Катедра, Економски факултет, УГД}
	
	% 1. Наслов
	\begin{frame}[plain]
		\titlepage
	\end{frame}
	
	% 2. Еден говорник. Ставете ја фотографијата до .tex датотеката; без неа се прикажуваат иницијалите.
	\UgdFrame{Отворање на настанот}{
		\UgdCard{0.72}{1.98}{14.56}{5.76}{ugdtint}
		\UgdAvatar{1.44}{2.58}{4.56}{govornik.jpg}{ИП}
		\node[ugdtxt,font=\ugdheadfont\bfseries\large,text width=8.3cm] at (6.72,-3.8) {Проф. д-р Име Презиме};
		\node[ugdtxt,text=ugdgrey,font=\normalsize,text width=8.3cm] at (6.72,-4.6) {Катедра};
	}
	
	% 3. Три лица
	\UgdFrame{Три лица}{
		\UgdPersonCard{0.72}{4.693}{13}{lice1.jpg}{ИП}{Име Презиме}{улога}
		\UgdPersonCard{5.653}{4.693}{13}{lice2.jpg}{ИП}{Име Презиме}{улога}
		\UgdPersonCard{10.586}{4.693}{13}{lice3.jpg}{ИП}{Име Презиме}{улога}
	}
	
	% 4. Четири лица
	\UgdFrame{Четири лица}{
		\UgdPersonCard{0.72}{3.46}{10}{lice1.jpg}{ИП}{Проф. д-р Име Презиме}{}
		\UgdPersonCard{4.42}{3.46}{10}{lice2.jpg}{ИП}{Проф. д-р Име Презиме}{}
		\UgdPersonCard{8.12}{3.46}{10}{lice3.jpg}{ИП}{Проф. д-р Име Презиме}{}
		\UgdPersonCard{11.82}{3.46}{10}{lice4.jpg}{ИП}{М-р Име Презиме}{}
	}
	
	% 5. Колони со голем број; последната е истакната
	\UgdFrame{Пет колони со голем број}{
		\foreach \i/\r/\topic/\n in {0/Дел 1/{Краток опис}/9,1/Дел 2/{Краток опис}/9,2/Дел 3/{Краток опис}/9,3/Дел 4/{Краток опис}/9}{
			\pgfmathsetmacro{\x}{0.72+\i*2.964}
			\UgdCard{\x}{1.98}{2.724}{4.2}{ugdtint}
			\node[ugdtxt,font=\normalsize\bfseries] at (\x+0.22,-2.22) {\r};
			\node[ugdtxt,text=ugdgrey,font=\scriptsize,text width=2.3cm] at (\x+0.22,-2.82) {\topic};
			\node[ugdtxt,text=ugdred,font=\ugdheadfont\bfseries\fontsize{28}{30}\selectfont] at (\x+0.22,-4.2) {\n};
			\node[ugdtxt,font=\small] at (\x+0.22,-5.4) {единици};
		}
		\UgdCard{12.576}{1.98}{2.724}{4.2}{ugdred}
		\node[ugdtxt,text=white,font=\normalsize\bfseries] at (12.796,-2.22) {Дел 5};
		\node[ugdtxt,text=white,font=\scriptsize] at (12.796,-2.82) {Истакнато};
		\node[ugdtxt,text=white,font=\ugdheadfont\bfseries\fontsize{28}{30}\selectfont] at (12.796,-4.2) {3};
		\node[ugdtxt,text=white,font=\small] at (12.796,-5.4) {единици};
		\node[ugdtxt,font=\normalsize] at (0.72,-6.72) {\textbf{Главен податок} \quad | \quad \textbf{Втор податок}};
	}
	
	% 6. Две панели: сива и црвена
	\UgdFrame{Две панели}{
		\UgdCard{0.72}{1.98}{7.08}{5.88}{ugdtint}
		\node[ugdtxt,font=\normalsize\bfseries] at (1.08,-2.22) {Наслов на панелот};
		\foreach \p in {1,...,9}{
			\pgfmathsetmacro{\x}{1.14+mod(\p-1,3)*1.14}
			\pgfmathsetmacro{\y}{3.06+floor((\p-1)/3)*1.14}
			\ifnum\p<4 \def\bg{ugdred}\def\fg{white}\else\def\bg{ugdline}\def\fg{ugdgrey}\fi
			\UgdDisc{\x}{\y}{0.96}{\bg}{\fg}{\p}
		}
		\node[anchor=west,inner sep=0pt,align=left] at (4.86,-3.54) {{\large\bfseries\color{ugdred}1--3}\\{\small истакнато}};
		\node[anchor=west,inner sep=0pt,align=left] at (4.86,-5.28) {{\large\bfseries\color{ugdgrey}4--9}\\{\small останато}};
		
		\UgdCard{8.196}{1.98}{7.08}{5.88}{ugdred}
		\node[ugdtxt,text=white,font=\normalsize\bfseries] at (8.556,-2.22) {Наслов на панелот};
		\foreach \i/\pl/\lab in {0/1/\textbf{Прва ставка},1/2/Втора ставка,2/3/Трета ставка}{
			\pgfmathsetmacro{\y}{3.06+\i*1.14}
			\UgdDisc{8.556}{\y}{0.96}{white}{ugdred}{\pl}
			\node[anchor=west,inner sep=0pt,text=white,font=\normalsize] at (9.78,-\y-0.48) {\lab};
		}
		\node[ugdtxt,text=white,font=\small,text width=6.3cm] at (8.556,-6.54) {Кратка забелешка во една или две реченици.};
	}
	
	% 7. Табела
	\UgdFrame{Табела}{
		\UgdCard{0.72}{1.98}{14.56}{0.82}{ugdred}
		\fill[ugdline] (0.72,-3.62) rectangle (15.28,-4.44);
		\foreach \y in {3.62,4.44,5.26,6.08}{\draw[ugdline] (0.72,-\y) -- (15.28,-\y);}
		\node[anchor=west,inner sep=0pt,text=white,font=\small\bfseries] at (0.96,-2.39) {Ставка};
		\foreach \x/\h in {5.5/Колона 1,8.5/Колона 2,11.5/Колона 3,13.86/Вкупно}{
			\node[inner sep=0pt,text=white,font=\small\bfseries] at (\x,-2.39) {\h};
		}
		\foreach \y/\t/\a/\b/\c/\s in {3.21/Ред 1/1/2/3/6,4.03/Ред 2/4/5/6/15,4.85/Ред 3/7/8/9/24,5.67/Ред 4/1/1/1/3}{
			\node[anchor=west,inner sep=0pt,font=\small\bfseries] at (0.96,-\y) {\t};
			\foreach \x/\v in {5.5/\a,8.5/\b,11.5/\c}{\node[inner sep=0pt,font=\small] at (\x,-\y) {\v};}
			\node[inner sep=0pt,font=\normalsize\bfseries] at (13.86,-\y) {\s};
		}
		\node[ugdtxt,font=\normalsize\bfseries,text width=14.5cm] at (0.72,-6.42) {Заклучок од табелата во една реченица.};
		\node[ugdtxt,text=ugdgrey,font=\scriptsize,text width=14.5cm] at (0.72,-7.14) {Извор или забелешка.};
	}
	
	% 8. Обичен слајд со точки и блок
	\begin{frame}[t]
		\frametitle{Обичен слајд}
		\begin{itemize}
			\item Прва точка со \alert{истакнат} збор.
			\item Втора точка.
		\end{itemize}
		\begin{block}{Наслов на блокот}
			Текст на блокот.
		\end{block}
	\end{frame}
	
	% 9. Крај
	\UgdClosingFrame{Ви благодарам на вниманието}{Кратка завршна порака.}
	
\end{document}
